% AUTORIDAD DE FUENTE:
% expose_VIII_body.tex, líneas 1607--1696; índices de escaneo 280--284;
% folios de fuente 269--273. No hay diferencias materiales entre las
% capas diplomática y corregida en este fragmento.

\medskip
\noindent\underline{Proposición 4.7.} \underline{Sean $P$ y $Q$ dos
grupos algebraicos lisos y conexos sobre el cuerpo $k$. Supongamos que
$P$ es una extensión de un esquema abeliano $A$ por un grupo algebraico
afín liso y conexo $L$ que admite una serie de composición cuyos
factores son isomorfos a un toro o a $\underline G_a$ (condición que se
cumple siempre si $k$ es perfecto}). \underline{Entonces el funtor considerado en
\sisi{1.3}{4.3} es una equivalencia de categorías; es decir, todo haz invertible
birrigidificado $\mathbb L$ sobre $P\times Q$ procede de una
biextensión de $(P,Q)$. Además, el funtor entre categorías rígidas}
\[
 \operatorname{BIEXT}(A,Q;\underline G_m)
 \longrightarrow
 \operatorname{BIEXT}(P,Q;\underline G_m)
\]
\underline{es una equivalencia de categorías; en otras palabras,
induce un isomorfismo de grupos}
\begin{equation}\tag{4.7.1}
 \Biext^1(A,Q;\underline G_m)
 \xrightarrow{\ \sim\ }
 \Biext^1(P,Q;\underline G_m).
\end{equation}

%%%%%%%%%%  scan idx 281  |  folio 270  %%%%%%%%%%

Para la primera afirmación, el descenso y 4.3 nos reducen de nuevo al
caso en que $k$ es separablemente cerrado. Procediendo como en la
demostración de 4.6, se deduce de 4.5 y de la hipótesis sobre $L$ que
$\mathbb L$ es isomorfo a la imagen inversa de un haz invertible
birrigidificado $\mathbb M$ sobre $A\times Q$. Para la primera
afirmación de 4.7 quedamos, pues, reducidos al caso $A=P$. El haz
invertible $\mathbb L$ define entonces un morfismo de esquemas
apuntados
\[
 f:Q\longrightarrow\underline{\Pic}_{P/k},
\]
que necesariamente se factoriza por el esquema abeliano dual $P'$ de
$P$. Sea $\mathbb L'$ la imagen inversa del haz de Weil sobre
$P\times P'$ mediante
\[
 \id_P\times f:P\times Q\longrightarrow P\times P'.
\]
Se trata de un haz invertible birrigidificado que, por definición de
$f$, solo difiere de $\mathbb L$ por la imagen inversa de un haz
invertible $\mathbb M$ sobre $P$. Como las restricciones de
$\mathbb L'$ y $\mathbb L$ a $P\times e_Q$ son triviales, se sigue que
$\mathbb M$ es trivial y, por tanto, que $\mathbb L\simeq\mathbb L'$.
Se señaló en VII 2.9.4 que las correspondencias divisoriales sobre el
producto de dos esquemas abelianos proceden de biextensiones. En
consecuencia, $\mathbb L'$, y por tanto $\mathbb L$, procede de una
biextensión.

Volvamos al caso general, en el que $P$ es una extensión de $A$ por
$L$. Queda por demostrar que (4.7.1) es un isomorfismo. Esto se deduce
de la sucesión exacta de los $\Biext^i$, para $i=0,1$, asociada a
\[
 0\longrightarrow L\longrightarrow P\longrightarrow A
 \longrightarrow0,
\]
y de las relaciones
\[
 \Biext^0(L,Q;\underline G_m)=0,
 \qquad
 \Biext^1(L,Q;\underline G_m)=0.
\]

%%%%%%%%%%  scan idx 282  |  folio 271  %%%%%%%%%%

La primera es un caso particular de la afirmación de rigidez de 4.3,
y la segunda es precisamente el primer caso de 4.6.

\medskip
\noindent\underline{Corolario 4.8.} \underline{Bajo las hipótesis de
4.7, supongamos que $Q$ es una extensión de un grupo algebraico $B$}
(\underline{necesariamente liso y conexo}) \underline{por un grupo algebraico afín liso y
conexo $M$. Entonces el funtor entre categorías rígidas}
\[
 \operatorname{BIEXT}(A,B;\underline G_m)
 \longrightarrow
 \operatorname{BIEXT}(P,Q;\underline G_m)
\]
\underline{es una equivalencia de categorías; en otras palabras,
induce un isomorfismo de grupos}
\begin{equation}\tag{4.8.1}
 \Biext^1(A,B;\underline G_m)
 \xrightarrow{\ \sim\ }
 \Biext^1(P,Q;\underline G_m).
\end{equation}

En virtud de (4.7.1), basta demostrar que
\[
 \Biext^1(A,B;\underline G_m)
 \xrightarrow{\ \sim\ }
 \Biext^1(A,Q;\underline G_m)
\]
es un isomorfismo. Como en la demostración de (4.7.1), esto se deduce
de la sucesión exacta de los $\Biext^i$ asociada a
\[
 0\longrightarrow M\longrightarrow Q\longrightarrow B
 \longrightarrow0,
\]
junto con las relaciones
\[
 \Biext^0(P,M;\underline G_m)=0,
 \qquad
 \Biext^1(P,M;\underline G_m)=0.
\]
En efecto, (4.7.1), aplicado con $M$ en lugar de $Q$, identifica estos
grupos con los grupos correspondientes que tienen $A$ en la primera
variable, tal como exige la sucesión exacta precedente. La segunda
anulación mostrada arriba es un caso particular del segundo caso de
4.6, tras intercambiar los dos factores.

%%%%%%%%%%  scan idx 283  |  folio 272  %%%%%%%%%%

\medskip
\noindent\underline{Observación 4.9.} Sean $P$ y $Q$ dos grupos
algebraicos lisos y conexos sobre un cuerpo \underline{perfecto} $k$.
Se sabe entonces [3] que cada uno es una extensión de un esquema
abeliano por un grupo liso, conexo y afín, y que este último es el
producto de un toro por una extensión iterada de copias de
$\underline G_a$ (véase la referencia dada en 4.6):
\[
 0\longrightarrow L\longrightarrow P\longrightarrow A
 \longrightarrow0,
 \qquad
 0\longrightarrow M\longrightarrow Q\longrightarrow B
 \longrightarrow0.
\]
Por tanto, 4.8 se aplica y muestra que la teoría de las biextensiones
de $(P,Q)$ por $\underline G_m$ se reduce a la teoría análoga (¡y mucho
más clásica!) para la pareja $(A,B)$ de variedades abelianas.

\medskip
\noindent\underline{Corolario 4.10.} \underline{Sean $P$ y $Q$ dos
grupos algebraicos lisos y conexos sobre $k$, sea $L$ un subgrupo
algebraico afín liso y conexo de $P$, y sea $E$ una biextensión de
$(P,Q)$ por $\underline G_m$. Por 2.2, $E$ induce un emparejamiento}
\[
 T_\ell(P)\times T_\ell(Q)\longrightarrow T_\ell(\underline G_m).
\]
\underline{Entonces el emparejamiento inducido}
\[
 T_\ell(L)\times T_\ell(Q)\longrightarrow T_\ell(\underline G_m)
\]
\underline{es nulo.}

En efecto, podemos suponer evidentemente que $k$ es algebraicamente
cerrado. Estamos entonces bajo las hipótesis de 4.6, que implican que
la restricción de $E$ a $(L,Q)$ es la biextensión trivial. La conclusión
afirmada se sigue porque el emparejamiento inducido en cuestión es el
definido por esta restricción.
